% ===== §13 =====
\section{Non-Action as the Highest Rationality}
\label{p25:sec:wuwei}

\noindent
The recovery closes with the moment that reaches furthest from optimization and gives the paper's positive claim its most demanding form: that the highest exercise of relational rationality is a kind of non-action, a knowing when not to intervene, not to force, not to optimize at all, and that this restraint, far from the absence of reason, is its most difficult accomplishment. The claim draws on the notion of \emph{wuwei}, non-action or effortless action, and completes the recovery by naming the competence that even the rationality of response must include: the competence of letting a generation arise rather than making it, of tending the conditions of a good rather than seizing the good directly \parencite{laozi2003daodejing}. It is the sharpest point of the whole positive account, because it is where relational rationality most flatly contradicts the imperative the optimizing frame cannot suspend, the imperative always to close the gap to the maximum.

\subsection*{The rationality of restraint}

The optimizing frame knows no restraint of this kind and cannot, for its imperative is always to act, to move toward the maximum, to close whatever distance remains between the present state and the best. Within that frame, inaction is intelligible only as a failure to optimize, a leaving of value on the table; there is no positive place in it for the deliberate refusal to intervene, since any refusal that forgoes an available improvement is, by the logic of the maximum, simply irrational. The rationality of relation has such a place, and it is central rather than marginal, because much of what is best in a shared life arises only when it is not forced and would be destroyed by the very intervention that sought to secure it. Trust that is demanded is not trust; spontaneity that is engineered is not spontaneous; a tenderness produced on purpose in order to produce it is something other than tenderness. For goods of this kind, and they are the characteristic goods of intimacy, the fitting exercise of reason is often to refrain: to make the space in which the good can arise of itself, and then to leave it be. And to know which goods are of this kind, and when the space has been made, and how to keep from the intervention that would spoil them, is a discernment as exacting as any the previous section described, and in one respect harder, since it must be exercised against the standing impulse to do something, to help, to secure, to optimize. This is the rationality of restraint, and it is not the abdication of reason but its refinement, the point at which reason grows wise enough to know the limits of its own doing and to recognize that some goods are available only to a hand that does not grasp.

\subsection*{Letting the good cycle arise}

The series has given this restraint a precise place before, and recalling it shows that the figure of non-action is not an exotic import brought in to decorate the ending but the name of a claim the argument has carried throughout. In the account of the good and the bad cycle, the good cycle was shown to be one whose positive accumulation is \emph{endogenous}, arising from the internal tension of the relation itself rather than pumped in by extraction from without; and non-action was there identified as the stance that neither extracts nor forces closure, that lets the endogenous accumulation occur of itself rather than manufacturing a false increase by drawing on some outside \parencite{paperix}. The same identification holds here and completes the argument of the present paper by joining its end to its middle. To optimize a relation is to intervene upon it, to act upon it as upon a system to be driven toward a maximum; and this intervention is exactly the forcing the good cycle cannot survive, since it substitutes an external drive for the internal generation and yields, at best, a manufactured increase that depletes what it draws upon, which is the very extraction the criticism of the third part described. The rationality of non-action is the refusal of this intervention: the tending of the conditions under which the relation generates value of itself, and the restraint from the optimizing action that would replace that self-generation with an extraction. What the maximizing frame reads as leaving value on the table is, seen rightly, the only way the value in question is ever had at all, since it is value that arises only where it is not seized, and the hand that reaches to maximize it closes on the diminished thing its reaching made.

\subsection*{The measure that is not the maximum}

The deepest form of the paper's positive thesis can now be stated, and it is fitting to end the recovery on it. Optimization takes its measure from the maximum: the standard against which every state is judged is the best attainable state, and reason's task is to close the distance to it. A rationality of the kind this part has described takes its measure elsewhere, from what a situation and an other are of themselves, from the grain of the thing and the movement it already has; and reason's task, on this measure, is to accord with that movement, to sustain and not force it, to tend and not seize, rather than to override it toward a maximum defined apart from it. This is the deepest sense in which optimization is neither the form of relational rationality nor the law of the world. The world of relation does not present a landscape with a summit that reason should climb. It presents a generation with a movement of its own that reason should discern and accord with; and the highest rationality is not the one that reaches the highest point but the one that accords most truly with what is there, that knows when to act and, harder, when not to, and finds in restraint not the failure of reason but its fulfillment. With this the recovery is complete. Relational rationality is response, attunement, sustaining, and, at its height, the non-action that lets the good arise of itself; and none of these is a maximization, and every one of them is reason.
